% Proposed editorial replacement sections for wcf_main.tex.
% These are integration-ready prose blocks. They intentionally omit
% internal version names, preregistration language, and intermediary hypotheses.

% Abstract
\begin{abstract}
Many interventions affect an outcome distribution rather than only its mean. This paper develops the Wasserstein Causal Forest (WCF), a distribution-valued learner for observational studies with binary treatment assignment and covariate-dependent treatment probabilities. WCF represents each conditional outcome law by particles on a monotone quantile grid, fits the particles with gradient boosting under an energy-score objective, uses a shared covariate partition for the treatment arms, and applies augmented inverse-propensity calibration to reported distributional functionals. We also give an optional two-part extension that combines a point mass with a positive-part particle law for outcomes with an exact-zero component. We evaluate the method in simulation studies covering placebo, location, distributional shape, and overlap changes, together with zero-inflated outcome laws. The estimands include treatment effects on distributional functionals and reference-adjusted TATE and TCATE contrasts, which measure whether treatment moves the outcome law toward a prespecified reference distribution. Across the reported regimes, the simulation evidence shows how a shared Wasserstein representation and functional calibration behave when treatment-arm support differs across prognostic covariates. The results identify the settings in which WCF is useful and make explicit the limitations imposed by finite particle representations, positivity, nuisance estimation, and the synthetic design.
\end{abstract}

% Introduction
\section{Introduction}

Average treatment effects reduce an intervention to a difference in means, although many policy and scientific questions concern changes in dispersion, tail behavior, participation, or other features of an outcome law. Conditional distributional treatment effects formalize these questions by comparing potential-outcome distributions or functionals of those distributions. Recent work studies conditional distributional effects through kernel conditional mean embeddings and U-statistic regression \citep{CATEGeneralization}, counterfactual mean embeddings \citep{Counterfactualmeanembeddings}, robust regression of distributional effect pseudo-outcomes \citep{GRFDistributionalCausalEffects}, kernel treatment-effect tests \citep{martineztaboada2023efficient}, and distributional random forests \citep{DRF-paper,naf2026causaldrf}. These approaches establish a methodological basis for asking causal questions about outcomes whose distributional structure is scientifically relevant.

We study the related problem in which the observed outcome for each unit is itself a distribution-valued object, represented here by a quantile function on a common grid. The Wasserstein geometry of quantile functions gives a natural metric for comparing such objects, while particle representations allow the conditional law of the quantile function to retain heterogeneity across units. The proposed Wasserstein Causal Forest combines this representation with a shared treatment-arm partition and a calibration step for scalar functionals of the law. The shared partition is intended for settings in which the treatment probability varies across covariates that also predict the potential-outcome law. This feature is distinct from the propensity score itself. The propensity score is the assignment probability $e(x)=\Pr(A=1\mid X=x)$ \citep{rosenbaum1983propensity}; covariate-dependent assignment becomes a source of confounding when those same covariates are prognostic and are not adequately adjusted for.

The paper defines two reference-adjusted estimands in addition to ordinary treatment-effect summaries. Let $Q^a$ denote the potential-outcome quantile vector for a unit, let $q_\star$ be a prespecified reference quantile vector, and let $d_W$ denote the grid Wasserstein distance. The reference-adjusted marginal effect is
\[
\mathrm{TATE}_{\mathrm{ref}}
=\E\left[d_W(Q^1,q_\star)-d_W(Q^0,q_\star)\right],
\]
so a negative value indicates that treatment moves the unit-level quantile law closer to the reference on average. Its conditional analogue replaces the outer expectation by conditioning on $X=x$ or on a moderator stratum, giving a reference-adjusted TCATE. These contrasts complement ordinary TATE and TCATE: they retain the direction of movement relative to a target distribution rather than only the distance between the two treatment arms.

The contribution is threefold. First, WCF provides a concrete Wasserstein and quantile representation for conditional laws of distribution-valued outcomes. Second, it combines a common treatment-arm partition with a standard augmented inverse-propensity score for calibration of arbitrary reported functionals. The calibration is doubly robust in the usual nuisance-model sense for each functional when either the propensity nuisance or the corresponding conditional functional regression is consistently estimated \citep{bang2005doubly,chernozhukov2018double}. Third, the paper introduces the reference-adjusted TATE and TCATE contrasts and gives an optional two-part assembly for laws with an exact-zero component. The latter is motivated by settings such as medical expenditure, where participation and positive intensity are naturally separated \citep{duan1983medical}, and by excess-zero count outcomes \citep{mullahy1986modified}.

We evaluate WCF in simulation studies designed to vary the treatment effect on location and distributional shape, the degree of overlap, the relationship between treatment assignment and prognostic covariates, and the presence of a point mass. The simulation design is stated directly in terms of the data-generating process. Income and wage distributions provide one interpretation of the quantile-valued outcome, but the same mathematical object can arise for regional expenditure distributions, health-care costs, benefit amounts, or any setting in which a unit is summarized by a distribution and an intervention acts on that distribution. The study therefore treats the income interpretation as an example of the design rather than as a claim about a completed empirical application.

% Discussion
\section{Discussion}

The simulation studies indicate that WCF can estimate distribution-valued treatment effects when treatment probabilities vary over prognostic covariates and the two treatment arms have different empirical support. The evidence is conditional on the data-generating processes, sample sizes, overlap ranges, particle budget, and reference distributions considered here. It supports a practical interpretation of the shared partition: both arms are represented in a common local geometry, while arm-specific contrasts are shrunk toward their pooled signal in leaves with limited arm information.

The calibration layer should be interpreted at the level at which it is defined. For a chosen functional $h$, the augmented inverse-propensity score combines an outcome regression for $h(q)$ with inverse-propensity residual correction. Under the standard causal identification assumptions and appropriate nuisance convergence, this score has the usual double-robust property \citep{bang2005doubly,chernozhukov2018double}. The property concerns the calibrated functional estimand. It does not imply that the finite particle law is doubly robust, that the shared tree partition is consistently selected under arbitrary misspecification, or that all Wasserstein functionals inherit the same finite-sample behavior. The propensity model used in the experiments is a logistic working model selected for simplicity; other binary classification learners can be substituted when their cross-fitted estimates satisfy the required positivity and convergence conditions.

The energy objective provides a proper distributional scoring principle before finite-particle smoothing and optimization are introduced. The energy score is strictly proper under its standard moment conditions \citep{gneiting2007proper}. In the implementation, the collision smoothing and fixed particle budget define a restricted computational objective. This choice makes the method tractable and preserves monotonicity through projection, but it also creates approximation error. Increasing the grid size and the number of particles can reduce discretization error at additional computational cost, and their sensitivity should be reported alongside the primary simulation results.

The two-part extension addresses a separate representation issue. A particle cloud with no explicit mixture gate can underrepresent a structural spike in the law of a distribution-valued outcome. When the observed law has a degenerate component at the zero quantile vector and a positive-part component, the extension estimates the mass probability and the positive-part law separately and then assembles them in law space. This construction is relevant to medical expenditures and other outcomes with structural zeros, but it is not required for outcomes whose distribution-valued support has no such degenerate component. The zero-inflation results should therefore be read as evidence about this optional extension under the stated mixture design.

Several limitations remain. The causal interpretation requires consistency, conditional exchangeability, and positivity, none of which can be verified from a simulation study or guaranteed in an observational application. Reference-adjusted effects depend on the choice of $Q_\star$, so the reference must be specified substantively before interpretation. Particle and quantile-grid approximations may omit features of a law that are important outside the chosen grid. Finally, the simulation study does not establish uniform dominance over distributional random forests or other conditional-law learners. It identifies performance patterns for the stated regimes and provides a basis for future comparisons using real distribution-valued outcomes.

% Conclusion
\section{Conclusion}

This paper develops a causal learner for outcomes represented by conditional distributions. WCF combines a Wasserstein quantile representation, particle boosting, a shared treatment-arm partition, and doubly robust calibration of distributional functionals. The reference-adjusted TATE and TCATE estimands provide an additional way to describe whether an intervention moves a conditional law toward a scientifically chosen target. Simulation studies show the behavior of the method under covariate-dependent assignment, varying overlap, distributional treatment effects, and exact-zero mixture components. The method is most naturally viewed as a distribution-valued extension of causal forest ideas whose validity depends on the usual causal assumptions, nuisance estimation, and finite representation choices. Applications with repeated regional, household, health, or expenditure distributions can assess whether these estimands capture policy-relevant changes beyond the mean.
